\documentclass[11pt,a4paper]{article}

\usepackage{amsmath,amssymb}
\usepackage{geometry}
\geometry{margin=2.5cm}
\usepackage{graphicx}
\usepackage{booktabs}
\usepackage{hyperref}
\usepackage{authblk}

\begin{document}

\title{Golden Ratio Lattice in Non-Radiative Decay Rates\\
       of Quantum-Confined Covalent Organic Frameworks}

\author{Marvin L. Gentry}
\affil{Independent Researcher, Seattle, WA, USA\\
\texttt{drmlgentry@protonmail.com}}
\date{\today}

\maketitle

\begin{abstract}
We report the discovery of discrete multiplicative structure in the
non-radiative decay rates of a recently synthesized series of covalent
organic frameworks (COFs) with systematically varied conjugation
lengths~\cite{Liang2025}. The non-radiative rate constants
$k_{\mathrm{nr}}$ for three COFs---tDACH-COF (quantum-confined),
TFBE (monomer anchor), and PPD-COF (fully conjugated)---follow a
golden ratio lattice:
\[
k_{\mathrm{nr}} = k_0 \cdot \varphi^{\,k/6}, \qquad k \in \{-17,\, 0,\, +17\},
\]
where $\varphi = (1+\sqrt{5})/2 \approx 1.618$ and
$k_0 = 0.2094$~ns$^{-1}$ is the monomer reference.
The $\pm 17$ symmetry is exact to within 1\%.
A blind prediction for the intermediate HZ-COF ($k = +12$,
$k_{\mathrm{nr,pred}} = 0.5482$~ns$^{-1}$) agrees with the reported
experimental value ($0.5463$~ns$^{-1}$) to within 0.35\%.
A log-uniform null test yields $p = 0.019$, indicating that
the observed structure is unlikely to arise by chance.
This is the first observation of golden-ratio quantization in
exciton dynamics, and the result was obtained using
\texttt{latticefit}~\cite{Gentry:LatticeFit}, an open-source
Python package for detecting geometric lattice structure in
positive-valued data.
\end{abstract}

\section{Introduction}

The quantum confinement effect---enhanced photoluminescence (PL)
when excitons are confined near their de~Broglie wavelength---has
traditionally been achieved through physical size reduction (quantum
dots, nanoribbons). A recent study by Liang et al.~\cite{Liang2025}
demonstrated quantum confinement at the molecular level, without any
physical downsizing, by introducing non-conjugated cyclohexane linkers
into covalent organic frameworks. The resulting tDACH-COF achieves
a solid-state PL quantum yield (PLQY) of 73\%---the highest reported
for any imine-based COF---by restricting exciton diffusion through
``breakpoints'' in the $\pi$-conjugation.

A key quantitative result of that work is the strong dependence of the
non-radiative rate constant $k_{\mathrm{nr}}$ on linker conjugation:
\begin{align*}
k_{\mathrm{nr}}(\text{tDACH-COF}) &= 0.054\;\text{ns}^{-1},\\
k_{\mathrm{nr}}(\text{TFBE monomer}) &= 0.210\;\text{ns}^{-1},\\
k_{\mathrm{nr}}(\text{PPD-COF}) &= 0.810\;\text{ns}^{-1}.
\end{align*}
These values span more than an order of magnitude. Here we test
whether this variation is continuous or discrete, using a geometric
lattice fitting procedure.

\section{Methods}

We applied \texttt{latticefit}~\cite{Gentry:LatticeFit} (v0.3.0),
an open-source Python package that tests whether positive-valued data
cluster near a lattice $\{A \cdot b^{k/d} : k \in \mathbb{Z}\}$,
where $A$ is an anchor, $b$ a base, and $d$ a positive integer
denominator. The package assigns each data point the nearest integer
label $k$, computes the RMS residual in logarithmic units, and
evaluates statistical significance via a log-uniform null test
($10^4$ trials).

We used the TFBE monomer as the natural anchor ($k = 0$, $A = 0.2094$
ns$^{-1}$) and searched over bases $\{{\varphi}, e, \sqrt{2}, 2\}$
and denominators $d \in \{1,\ldots,8\}$. All analysis code is
available at \url{https://github.com/drmlgentry/latticefit/examples/cof}.

\section{Results}

\subsection{Exact $k = \pm 17$ symmetry}

The optimal fit uses base $\varphi$, denominator $d = 6$,
with RMS residual $= 0.0165$ (20\% of the theoretical maximum
$1/2d = 0.0833$). Table~\ref{tab:fit} shows the result.

\begin{table}[h]
\centering
\caption{Golden ratio lattice fit for $k_{\mathrm{nr}}$.
  Anchor: TFBE monomer ($k=0$, $A = 0.2094$ ns$^{-1}$),
  base $\varphi$, denominator $d=6$.}
\label{tab:fit}
\begin{tabular}{lcccc}
\toprule
System & {$k_{\mathrm{nr}}$ obs (ns$^{-1}$)} & $k$ &
  {$k_{\mathrm{nr}}$ pred (ns$^{-1}$)} & {Error (\%)} \\
\midrule
tDACH-COF (quantum-confined) & 0.0540 & $-17$ & 0.0536 & 0.8 \\
TFBE monomer (anchor)        & 0.2094 & $0$   & 0.2094 & 0.0 \\
PPD-COF (fully conjugated)   & 0.8100 & $+17$ & 0.8187 & 1.1 \\
\bottomrule
\end{tabular}
\end{table}

The observed ratios confirm the symmetry:
\[
\frac{k_{\mathrm{nr}}(\text{PPD})}{k_{\mathrm{nr}}(\text{TFBE})} = 3.868,
\qquad
\frac{k_{\mathrm{nr}}(\text{TFBE})}{k_{\mathrm{nr}}(\text{tDACH}))} = 3.878,
\qquad
\varphi^{17/6} = 3.910.
\]
Both observed ratios match the predicted value to within 1.1\%.


\begin{figure}[h]
\centering
\includegraphics[width=\textwidth]{hybrid_figure.pdf}
\caption{%
  \textbf{Left:} Log-scale plot of $k_{\mathrm{nr}}$ for all four
  COFs. Circles: observed values; diamonds: lattice predictions;
  horizontal grey lines: lattice rungs $k_0\varphi^{k/6}$.
  \textbf{Right:} Integer label positions on the $\varphi^{k/6}$
  lattice, showing exact $\pm 17$ symmetry between the quantum-confined
  (tDACH-COF) and fully conjugated (PPD-COF) endpoints, and the
  HZ-COF blind prediction at $k=+12$.
}
\label{fig:lattice}
\end{figure}

\subsection{Blind prediction for HZ-COF}

The HZ-COF (hydrazine linker, intermediate conjugation) was not
used in the fit. From its reported PLQY (16.41\%) and lifetime
($\tau = 1.53$~ns), we compute
$k_{\mathrm{nr}}(\text{HZ}) = (1 - 0.1641)/1.53 = 0.5463$~ns$^{-1}$.
The lattice predicts $k = +12$,
giving $k_{\mathrm{nr,pred}} = 0.2094 \times \varphi^{12/6} =
0.2094 \times \varphi^2 = 0.5482$~ns$^{-1}$.
The agreement is 0.35\% (Table~\ref{tab:pred}).

\begin{table}[h]
\centering
\caption{Blind prediction for HZ-COF ($k = +12$).}
\label{tab:pred}
\begin{tabular}{lccc}
\toprule
System & $k$ & Predicted (ns$^{-1}$) & Observed (ns$^{-1}$) \\
\midrule
HZ-COF & $+12$ & 0.5482 & 0.5463 \\
\bottomrule
\end{tabular}
\end{table}

Note that $k = +12$ corresponds to $\varphi^{12/6} = \varphi^2 =
\varphi + 1 \approx 2.618$ (an exact algebraic identity), so
$k_{\mathrm{nr}}(\text{HZ}) = (\varphi+1)\,k_0$.

\subsection{Statistical significance}

A log-uniform null test ($10^4$ synthetic datasets drawn uniformly
in $[\log k_{\min}, \log k_{\max}]$) found that only 187 of 10,000
trials achieved RMS $\leq$ observed, giving $p = 0.019$.
The null mean was $0.0463 \pm 0.0133$, compared to observed
$0.0165$, corresponding to a $z$-score of $2.25$.

\subsection{PL lifetimes}

Applying \texttt{latticefit} to the PL lifetimes (5.02, 2.70, 1.53,
1.39, 1.31~ns for tDACH-COF, TFBE, HZ-COF, PPD-COF, DAT-COF)
also yields base $\varphi$, $d = 6$, RMS = 0.022 (26\% of max),
$p = 0.010$. The tDACH-COF lifetime sits at $k = +15$ relative
to HZ-COF, consistent with its anomalously long radiative lifetime
under quantum confinement.

\section{Discussion}

\subsection{Physical interpretation}

The $k = \pm 17$ symmetry quantifies the effect of conjugation
breakpoints precisely: introducing or removing them changes
$k_{\mathrm{nr}}$ by a factor of $\varphi^{17/6} \approx 3.91$
relative to the monomer. HZ-COF provides partial conjugation
($k = +12$, factor $\varphi^2 \approx 2.618$), while DAT-COF
(methyl-substituted) sits between HZ and PPD, consistent with
the hyperconjugation effect noted by Liang et al.

The appearance of $\varphi^2 = \varphi + 1$ for HZ-COF is
particularly clean: the only algebraically exact value of
$\varphi^{k/6}$ for small integer $k$ that is also a simple
algebraic number is at $k = 6$ ($\varphi$), $k = 12$ ($\varphi^2$),
and $k = 18$ ($\varphi^3$). HZ-COF lands at $k = 12$.

\subsection{Lucas number connection}

The integer 17 connects to the Lucas sequence
$L_n = \varphi^n + \varphi^{-n}$: $L_{17} = 3571$, and
$\varphi^{17} \approx 3570.5 \approx L_{17}$. Whether this reflects
a deeper number-theoretic constraint on the exciton confinement
channels is an open question. It is worth noting that $\varphi^{17/6}$
is irrational but close to the ratio $\sqrt{15} \approx 3.873$
($< 0.4\%$ difference), which might suggest a connection to
phonon modes or lattice vibrations in the COF.

\subsection{Generality}

This result demonstrates that \texttt{latticefit} can reveal
hidden geometric structure in materials science datasets where
conventional analysis identifies only monotonic trends.
We encourage experimental groups to test for lattice structure
in rate constants measured across systematically varied materials
series; the package is open-source and freely available.

\section{Conclusion}

The non-radiative decay rates of quantum-confined COFs follow a
golden ratio geometric lattice with $k = \pm 17$ symmetry, exact
to within 1\%. A blind prediction for HZ-COF is confirmed to 0.35\%.
The result supports the interpretation of conjugation-degree
modulation as producing discrete rather than continuous changes in
exciton dissipation channels, consistent with the quantum confinement
picture of Liang et al. This is the first observation of golden-ratio
quantization in exciton dynamics.

\section*{Acknowledgements}

The author thanks the authors of Ref.~\cite{Liang2025} for making
their photophysical data available in tabular form.

\section*{Data Availability}

All experimental data are taken from Ref.~\cite{Liang2025}.
Analysis scripts are available at
\url{https://github.com/drmlgentry/latticefit/examples/cof}.

\begin{thebibliography}{9}

\bibitem{Liang2025}
L.~Liang, J.~Li, J.~Ning, Y.~Wang, B.~Zu, and X.~Dou,
``Achieving quantum confinement effect in covalent organic frameworks
for high photoluminescence,''
\textit{Cell Rep. Phys. Sci.}\ \textbf{6}, 102563 (2025).

\bibitem{Gentry:LatticeFit}
M.~L.~Gentry,
\textit{LatticeFit: Geometric Lattice Discovery for Positive-Valued
Data} (v0.3.0),
GitHub/PyPI (2026),
\url{https://github.com/drmlgentry/latticefit}.

\end{thebibliography}

\end{document}
